% notes/latex/demo.tex
% Tier 1 LaTeX demo: exercises every block environment from the block-notes
% template exactly once, plus a pgfplots Ellingham chart, a TikZ free-body
% diagram, a CircuiTikZ circuit, an mhchem reaction, and a minted Python block.
%
% Build (needs -shell-escape for minted + biber for the bibliography):
%   latexmk -pdf -shell-escape demo.tex
% from notes/latex/, or from the repo root:
%   latexmk -pdf -shell-escape -cd notes/latex/demo.tex
\documentclass[11pt,a4paper]{article}

\usepackage[utf8]{inputenc}
\usepackage[T1]{fontenc}
\usepackage[margin=25mm]{geometry}
\usepackage{graphicx}

% Load the shared block-notes template (packages + tcolorbox environments).
% Path is relative to this file; latexmk runs with -cd so this resolves.
% templates/block-notes.tex
% Block-style engineering / materials-science notes template.
%
% This file is a *preamble fragment*: it loads the full LaTeX stack and defines
% seven reusable tcolorbox environments. Documents \input this file inside their
% own preamble (see notes/latex/demo.tex), then set \bibliography data and call
% \begin{document}. Build with latexmk + -shell-escape (minted needs Pygments):
%
%   latexmk -pdf -shell-escape notes/latex/demo.tex
%
% Requires (Tier 1 apt): texlive-latex-extra texlive-science texlive-pictures
% texlive-publishers texlive-bibtex-extra latexmk biber; python-pygments for
% minted. Bibliography backend: biber (biblatex).

% --- math, units, chemistry -------------------------------------------------
\usepackage{amsmath}
\usepackage{amssymb}
\usepackage{siunitx}
\usepackage[version=4]{mhchem}
\usepackage{chemfig}

% --- graphics / diagrams ----------------------------------------------------
\usepackage{tikz}
\usepackage{pgfplots}
\pgfplotsset{compat=1.18}
\usetikzlibrary{arrows.meta,positioning,calc}
\usepackage{circuitikz}

% --- code listings ----------------------------------------------------------
\usepackage{minted}

% --- colored block environments ---------------------------------------------
\usepackage[breakable,skins,minted]{tcolorbox}

% --- bibliography / links ---------------------------------------------------
\usepackage[backend=biber,style=numeric,sorting=none]{biblatex}
\usepackage{hyperref}
\hypersetup{colorlinks=true,linkcolor=blue!55!black,citecolor=teal!60!black,urlcolor=magenta!60!black}

% ---------------------------------------------------------------------------
% Seven reusable block environments. Each is breakable so long notes flow
% across pages. Titles are fixed so notes stay visually consistent.
% ---------------------------------------------------------------------------

% 1. theorybox - statements of theory / definitions.
\newtcolorbox{theorybox}[1][]{%
  breakable, enhanced, colback=blue!5!white, colframe=blue!60!black,
  fonttitle=\bfseries, title=Theory, #1}

% 2. derivationbox - step-by-step derivations.
\newtcolorbox{derivationbox}[1][]{%
  breakable, enhanced, colback=purple!5!white, colframe=purple!55!black,
  fonttitle=\bfseries, title=Derivation, #1}

% 3. diagrambox - figures / diagrams (TikZ, circuits, FBDs).
\newtcolorbox{diagrambox}[1][]{%
  breakable, enhanced, colback=teal!6!white, colframe=teal!55!black,
  fonttitle=\bfseries, title=Diagram, #1}

% 4. codebox - minted source listings (uses the tcolorbox minted integration).
\newtcolorbox{codebox}[1][]{%
  breakable, enhanced, colback=gray!8!white, colframe=gray!55!black,
  fonttitle=\bfseries, title=Code, #1}

% 5. examplebox - worked examples.
\newtcolorbox{examplebox}[1][]{%
  breakable, enhanced, colback=green!6!white, colframe=green!50!black,
  fonttitle=\bfseries, title=Example, #1}

% 6. warningbox - cautions / pitfalls.
\newtcolorbox{warningbox}[1][]{%
  breakable, enhanced, colback=red!6!white, colframe=red!60!black,
  fonttitle=\bfseries, title=Warning, #1}

% 7. datatable - tabular data / measured values.
\newtcolorbox{datatable}[1][]{%
  breakable, enhanced, colback=orange!6!white, colframe=orange!65!black,
  fonttitle=\bfseries, title=Data, #1}


% Bibliography database (biblatex/biber).
\addbibresource{../../refs/library.bib}

\title{Tier 1 Demo: Notes, Diagrams, and Plotting}
\author{Scientific-Computing Workspace}
\date{\today}

\begin{document}
\maketitle

% --- 1. theorybox ----------------------------------------------------------
\begin{theorybox}
The Gibbs free energy of a reaction, \(\Delta G = \Delta H - T\,\Delta S\),
governs oxide stability. An Ellingham diagram plots \(\Delta G\) versus
temperature \(T\) for metal oxidation, as discussed by \textcite{ellingham1944}.
\end{theorybox}

% --- 2. derivationbox ------------------------------------------------------
\begin{derivationbox}
For carbothermic reduction the relevant lines cross where carbon monoxide
formation becomes favorable:
\begin{align}
  \Delta G_{\text{ox}} &= \Delta H_{\text{ox}} - T\,\Delta S_{\text{ox}}, \\
  \Delta G_{\ce{CO}}   &= \Delta H_{\ce{CO}} - T\,\Delta S_{\ce{CO}}, \\
  T_{\text{eq}}        &= \frac{\Delta H_{\text{ox}} - \Delta H_{\ce{CO}}}
                              {\Delta S_{\text{ox}} - \Delta S_{\ce{CO}}} .
\end{align}
The overall reduction of hematite by carbon monoxide is
\ce{Fe2O3 + 3CO -> 2Fe + 3CO2}.
\end{derivationbox}

% --- 3. diagrambox: Ellingham pgfplots + TikZ FBD + CircuiTikZ -------------
\begin{diagrambox}
Ellingham-style plot of \(\Delta G\) against temperature (schematic):

\begin{center}
\begin{tikzpicture}
\begin{axis}[
    width=0.8\linewidth, height=5cm,
    xlabel={Temperature \(T\) (\si{\kelvin})},
    ylabel={\(\Delta G\) (\si{\kilo\joule\per\mole})},
    legend pos=south west, grid=both, domain=300:1800, samples=2,
  ]
  \addplot[blue, thick]  {-560 + 0.17*x};  \addlegendentry{\ce{2Fe + 3/2 O2 -> Fe2O3}}
  \addplot[red,  thick]  {-230 - 0.17*x};  \addlegendentry{\ce{2CO -> 2C + O2}}
\end{axis}
\end{tikzpicture}
\end{center}

Free-body diagram (TikZ) of a block on an incline:

\begin{center}
\begin{tikzpicture}[>=Stealth, scale=1.1]
  \draw[thick] (0,0) -- (4,0);
  \filldraw[fill=gray!20] (1.6,0.05) rectangle (2.4,0.75);
  \coordinate (c) at (2.0,0.4);
  \draw[->, thick, blue]  (c) -- ++(0,1.2)  node[above] {\(N\)};
  \draw[->, thick, red]   (c) -- ++(0,-1.2) node[below] {\(mg\)};
  \draw[->, thick, teal]  (c) -- ++(1.2,0)  node[right] {\(F\)};
  \draw[->, thick, orange](c) -- ++(-1.0,0) node[left]  {\(f\)};
\end{tikzpicture}
\end{center}

A simple RC circuit (CircuiTikZ):

\begin{center}
\begin{circuitikz}[scale=1.0]
  \draw (0,0) to[V, l=\(V_s\)] (0,2)
        to[R, l=\(R\)] (2,2)
        to[C, l=\(C\)] (2,0) -- (0,0);
\end{circuitikz}
\end{center}
\end{diagrambox}

% --- 4. codebox: minted Python ---------------------------------------------
\begin{codebox}
\begin{minted}{python}
import numpy as np

def delta_g(T, dH, dS):
    """Gibbs free energy line: dG = dH - T*dS (kJ/mol)."""
    return dH - T * dS

T = np.linspace(300, 1800, 16)
print(delta_g(T, dH=-560.0, dS=0.17))
\end{minted}
\end{codebox}

% --- 5. examplebox ---------------------------------------------------------
\begin{examplebox}
At \(T = \SI{1200}{\kelvin}\), using \(\Delta H = \SI{-560}{\kilo\joule\per\mole}\)
and \(\Delta S = \SI{0.17}{\kilo\joule\per\mole\per\kelvin}\):
\(\Delta G = -560 - 1200 \times 0.17 = \SI{-356}{\kilo\joule\per\mole}\).
\end{examplebox}

% --- 6. warningbox ---------------------------------------------------------
\begin{warningbox}
Ellingham lines assume all species are at unit activity and ignore kinetics;
a thermodynamically favorable reduction may still be kinetically slow.
\end{warningbox}

% --- 7. datatable ----------------------------------------------------------
\begin{datatable}
\begin{tabular}{lrr}
\hline
Species & \(\Delta H\) (\si{\kilo\joule\per\mole}) & \(\Delta S\) (\si{\joule\per\mole\per\kelvin}) \\
\hline
\ce{Fe2O3} & \num{-824} & \num{87}  \\
\ce{CO}    & \num{-111} & \num{198} \\
\ce{CO2}   & \num{-394} & \num{214} \\
\hline
\end{tabular}
\end{datatable}

\printbibliography

\end{document}
